% ============================================================
%  Experiment 3 --- Introduction to Altera Quartus II Software
% ============================================================
\experimentheader{3}{Introduction to Quartus II}{Introduction to Altera Quartus II Software}{FPGA Lab}

\begin{objectives}
\begin{enumerate}
  \item To use schematic entry as a design method in Altera Quartus II.
  \item To use Verilog as a design entry method in Altera Quartus II.
  \item To combine schematic and Verilog as a single design entry in Altera Quartus II.
  \item To simulate the circuit using Altera Quartus II.
  \item To download the design onto an Altera DE2 board.
\end{enumerate}
\end{objectives}

\begin{equipment}
\begin{itemize}
  \item Altera Quartus II software
  \item Altera DE2 board
\end{itemize}
\end{equipment}

\begin{procedure}

\textbf{Activity 1 --- Schematic entry}\vspace{2pt}

\textit{Objective: use schematic entry to design a 4-to-1 multiplexer in Quartus II.}

\begin{enumerate}
  \item Open Quartus II and start the New Project Wizard.
  \begin{itemize}
    \item Working directory: \texttt{fpga}; project name: \texttt{activity1}; top-level entity: \texttt{mux\_sch}.
    \item Device family: Cyclone II, device \texttt{EP2C35F672C6N}.
    \item Accept the default EDA tool settings (synthesis, simulation, timing analysis) and finish.
  \end{itemize}
  \item Create a new Block Diagram/Schematic file.
  \item Place three 3-input AND gates and one 4-input OR gate from the primitives library, and wire them as shown in Figure~\ref{fig:exp3-mux-schematic} using rubberbanding to connect components.
  \item Add and name the input pins (\texttt{i0}--\texttt{i3}, \texttt{s0}, \texttt{s1}) and the output pin (\texttt{out}).
  \item Save as \texttt{mux\_sch.bdf} and start compilation.
  \item Create a Vector Waveform File to simulate the design: insert input nodes \texttt{i0}--\texttt{i3}, \texttt{s0}, \texttt{s1}, and output node \texttt{out}.
  \item Assign logic values to the inputs, save as \texttt{mux\_sch.vwf}, then run Compilation and Simulation.
  \item Record the outputs, then repeat for each input combination in Table~\ref{tab:exp3-truth}, saving the waveform file each time.
  \item Print and submit the schematic diagram and simulation result.
\end{enumerate}

\begin{boxedfigure}
\centering
\begin{tikzpicture}[node distance=6mm and 10mm]
  \node[io] (i0) {i0};
  \node[io, below=4mm of i0] (i1) {i1};
  \node[io, below=4mm of i1] (i2) {i2};
  \node[io, below=4mm of i2] (i3) {i3};
  \node[io, below=6mm of i3] (s0) {s0};
  \node[io, below=4mm of s0] (s1) {s1};
  \node[block, right=26mm of i1, minimum height=3.6cm, minimum width=2.4cm, yshift=-6mm] (mux) {4-to-1\\MUX};
  \node[io, right=14mm of mux] (out) {out};
  \foreach \sig in {i0,i1,i2,i3} { \draw[arr] (\sig) -- (\sig -| mux.west); }
  \draw[arr] (s0) -| ($(mux.south)+(-0.5,0)$);
  \draw[arr] (s1) -| ($(mux.south)+(0.5,0)$);
  \draw[arr] (mux) -- (out);
\end{tikzpicture}
\captionof{figure}{Schematic-entry 4-to-1 multiplexer (\texttt{mux\_sch}).}
\label{fig:exp3-mux-schematic}
\end{boxedfigure}

\begin{boxedtable}
\centering
\begin{tabular}{cccccc}
\toprule
S0 & S1 & I0 & I1 & I2 & I3 \\
\midrule
0 & 0 & 1 & 1 & 0 & 0 \\
0 & 1 & 0 & 1 & 1 & 0 \\
1 & 0 & 1 & 1 & 0 & 1 \\
1 & 1 & 1 & 1 & 0 & 0 \\
\bottomrule
\end{tabular}
\captionof{table}{Input combinations to simulate.}
\label{tab:exp3-truth}
\end{boxedtable}

\textbf{Activity 2 --- Verilog entry}\vspace{2pt}

\textit{Objective: use Verilog code as a design entry method in Quartus II.}

\begin{enumerate}
  \item Start a new project (\texttt{Activity2}, top level \texttt{mux\_ver}) with the same device and EDA tool settings as Activity 1.
  \item Create a new Verilog file and enter the code in Listing~\ref{lst:exp3-mux-ver}.
  \item Save as \texttt{mux\_ver}, then compile the design.
  \item Create a Vector Waveform File; insert input nodes \texttt{i0}--\texttt{i3}, \texttt{s0}, \texttt{s1}, and output node \texttt{out}.
  \item Assign logic values matching Table~\ref{tab:exp3-truth}, save as \texttt{mux\_ver.vwf}, then start simulation.
  \item Inspect the result in the Simulation Report --- Simulation Waveform.
  \item Change the input values (saving before each run) and re-simulate.
  \item View the synthesized schematic via \textbf{Tools} $\to$ \textbf{Netlist Viewer} $\to$ \textbf{RTL Viewer}, shown in Figure~\ref{fig:exp3-rtl}.
  \item Print and submit the Verilog code and simulation result.
\end{enumerate}

\begin{codepanel}{mux\_ver.v}{verilogstyle}
\begin{lstlisting}
module mux_ver(out, i0, i1, i2, i3, s1, s0);
  output out;
  input i0, i1, i2, i3;
  input s1, s0;
  reg out;

  always @(s1 or s0 or i0 or i1 or i2 or i3)
  begin
    case ({s1, s0})
      2'b00: out = i0;
      2'b01: out = i1;
      2'b10: out = i2;
      2'b11: out = i3;
      default: out = 1'bx;
    endcase
  end
endmodule
\end{lstlisting}
\end{codepanel}
\vspace{-4pt}
\begin{center}\small\color{slategray}Listing \refstepcounter{equation}\theequation\label{lst:exp3-mux-ver}: \texttt{mux\_ver.v} --- Verilog 4-to-1 multiplexer.\end{center}

\begin{boxedfigure}
\centering
\begin{tikzpicture}
  \node[block, minimum width=1.6cm, minimum height=2.6cm, trapezium, trapezium left angle=70, trapezium right angle=110, shape border rotate=0] (mux) {};
  \node[above=1mm of mux, font=\small\plexmedium] {Mux0};
  \node[io, left=12mm of mux, yshift=8mm] (s1) {s1};
  \node[io, below=3mm of s1] (s0) {s0};
  \node[io, below=5mm of s0] (i3) {i3};
  \node[io, below=3mm of i3] (i2) {i2};
  \node[io, below=3mm of i2] (i1) {i1};
  \node[io, below=3mm of i1] (i0) {i0};
  \node[io, right=14mm of mux] (out) {out};
  \draw[arr] (s1) -- ($(mux.west)+(0,0.9)$) node[midway, above, font=\tiny] {SEL[1..0]};
  \draw[arr] (s0) -- ($(mux.west)+(0,0.55)$);
  \draw[arr] (i3) -- ($(mux.west)+(0,-0.1)$) node[midway, above, font=\tiny] {DATA[3..0]};
  \draw[arr] (i2) -- ($(mux.west)+(0,-0.35)$);
  \draw[arr] (i1) -- ($(mux.west)+(0,-0.6)$);
  \draw[arr] (i0) -- ($(mux.west)+(0,-0.85)$);
  \draw[arr] (mux.east) -- (out) node[midway, above, font=\tiny] {out};
\end{tikzpicture}
\captionof{figure}{RTL Viewer output for \texttt{mux\_ver} (\texttt{Mux0}).}
\label{fig:exp3-rtl}
\end{boxedfigure}

\textbf{Activity 3 --- Combining schematic and Verilog entry}\vspace{2pt}

\textit{Objective: combine a Verilog module with schematic entry in a single design.}

\begin{enumerate}
  \item Start a new project (\texttt{activity3}, top level \texttt{mux\_schver}) with the same device and EDA tool settings.
  \item Create a new Verilog HDL file with the code in Listing~\ref{lst:exp3-mux-vrl}; save as \texttt{mux\_vrl}.
  \item Generate a reusable symbol from the file via \textbf{File} $\to$ \textbf{Create/Update} $\to$ \textbf{Create Symbol Files for Current File}. The resulting block is shown in Figure~\ref{fig:exp3-mux-vrl-symbol}.
  \item Create a new Block Diagram/Schematic file, place the \texttt{mux\_vrl} symbol from the project library, and connect its inputs and output to input/output pads as shown in Figure~\ref{fig:exp3-mux-schver}.
  \item Save as \texttt{mux\_schver.bdf} and start compilation.
  \item Create a Vector Waveform File, insert the same nodes as before, assign logic values, save as \texttt{mux\_schver.vwf}, and run Compilation and Simulation.
  \item Record the outputs for the combinations in Table~\ref{tab:exp3-truth}, saving the waveform file each time.
  \item Print and submit the schematic diagram, Verilog code, and simulation result.
\end{enumerate}

\begin{codepanel}{mux\_vrl.v}{verilogstyle}
\begin{lstlisting}
module mux_vrl(out, i0, i1, i2, i3, s1, s0);
  output out;
  input i0, i1, i2, i3;
  input s1, s0;
  reg out;

  always @(s1 or s0 or i0 or i1 or i2 or i3)
  begin
    case ({s1, s0})
      2'b00: out = i0;
      2'b01: out = i1;
      2'b10: out = i2;
      2'b11: out = i3;
      default: out = 1'bx;
    endcase
  end
endmodule
\end{lstlisting}
\end{codepanel}
\vspace{-4pt}
\begin{center}\small\color{slategray}Listing \refstepcounter{equation}\theequation\label{lst:exp3-mux-vrl}: \texttt{mux\_vrl.v} --- Verilog module to be reused as a schematic symbol.\end{center}

\begin{boxedfigure}
\centering
\begin{tikzpicture}
  \node[block, minimum width=2.4cm, minimum height=3.2cm] (sym) {};
  \node[above=1mm of sym, font=\small\plexmedium] {mux\_vrl};
  \node[font=\scriptsize, anchor=west] at ($(sym.north west)+(0.15,-0.35)$) {i0};
  \node[font=\scriptsize, anchor=west] at ($(sym.north west)+(0.15,-0.75)$) {i1};
  \node[font=\scriptsize, anchor=west] at ($(sym.north west)+(0.15,-1.15)$) {i2};
  \node[font=\scriptsize, anchor=west] at ($(sym.north west)+(0.15,-1.55)$) {i3};
  \node[font=\scriptsize, anchor=west] at ($(sym.north west)+(0.15,-1.95)$) {s1};
  \node[font=\scriptsize, anchor=west] at ($(sym.north west)+(0.15,-2.35)$) {s0};
  \node[font=\scriptsize, anchor=east] at ($(sym.north east)+(-0.15,-0.35)$) {out};
\end{tikzpicture}
\captionof{figure}{Auto-generated \texttt{mux\_vrl} symbol block.}
\label{fig:exp3-mux-vrl-symbol}
\end{boxedfigure}

\begin{boxedfigure}
\centering
\begin{tikzpicture}[node distance=4mm and 14mm]
  \node[io] (i0) {i0};
  \node[io, below=3mm of i0] (i1) {i1};
  \node[io, below=3mm of i1] (i2) {i2};
  \node[io, below=3mm of i2] (i3) {i3};
  \node[io, below=3mm of i3] (s1) {s1};
  \node[io, below=3mm of s1] (s0) {s0};
  \node[block, right=of i2, minimum height=3.4cm, minimum width=2.2cm, yshift=-9mm] (mux) {mux\_vrl};
  \node[io, right=14mm of mux] (out) {out};
  \foreach \sig in {i0,i1,i2,i3,s1,s0} { \draw[arr] (\sig) -- (\sig -| mux.west); }
  \draw[arr] (mux) -- (out);
\end{tikzpicture}
\captionof{figure}{Combined schematic \texttt{mux\_schver}: the Verilog-derived \texttt{mux\_vrl} block wired to I/O pads.}
\label{fig:exp3-mux-schver}
\end{boxedfigure}

\textbf{Activity 4 --- Timing analysis}\vspace{2pt}

\textit{Objective: analyse the maximum operating speed of a design.}

\begin{enumerate}
  \item Start a new project (\texttt{Activity4}, top level \texttt{counter}) with the same device and EDA tool settings.
  \item Create a new Verilog file with the code in Listing~\ref{lst:exp3-counter}; save as \texttt{counter} and compile.
  \item Create a Vector Waveform File; insert \texttt{clock} and \texttt{clear} as inputs and \texttt{Q} as output.
  \item Force \texttt{clear} low, assign a clock signal to \texttt{clock}, save as \texttt{counter.vwf}, then start simulation.
  \item Inspect the result in the Simulation Report --- Simulation Waveform.
  \item Run \textbf{Processing} $\to$ \textbf{Start} $\to$ \textbf{Start Classic Timing Analyzer}.
  \item Read the maximum clock frequency from the Timing Analyzer Summary (Table~\ref{tab:exp3-timing}).
  \item Print and submit the Verilog code and simulation result.
\end{enumerate}

\begin{codepanel}{counter.v}{verilogstyle}
\begin{lstlisting}
module counter(Q, clock, clear);
  output [3:0] Q;
  input clock, clear;
  reg [3:0] Q;

  always @(posedge clear or negedge clock)
  begin
    if (clear)
      Q = 4'd0;
    else
      Q = (Q + 1) % 16;
  end
endmodule
\end{lstlisting}
\end{codepanel}
\vspace{-4pt}
\begin{center}\small\color{slategray}Listing \refstepcounter{equation}\theequation\label{lst:exp3-counter}: \texttt{counter.v} --- 4-bit up counter with asynchronous clear.\end{center}

\begin{boxedtable}
\centering
\begin{tabular}{clll}
\toprule
\# & Type & Actual time & Failed paths \\
\midrule
1 & Worst-case $t_{co}$ (clock $\to$ Q[3]) & 7.410\,ns & 0 \\
2 & Clock setup (\texttt{clock}) & 233.64\,MHz (4.280\,ns period) & 0 \\
3 & Total failed paths & --- & 0 \\
\bottomrule
\end{tabular}
\captionof{table}{Timing Analyzer Summary for \texttt{counter}.}
\label{tab:exp3-timing}
\end{boxedtable}

\textbf{Activity 5 --- Downloading to the DE2 board}\vspace{2pt}

\textit{Objective: download a Verilog design onto the Altera DE2 board and drive the seven-segment display.}

\begin{enumerate}
  \item Start a new project (\texttt{activity5}, top level \texttt{counta}) with the same device and EDA tool settings.
  \item Create a new Verilog HDL file with the code in Listing~\ref{lst:exp3-counta}; save as \texttt{counta}.
  \item Compile and simulate the design.
  \item Assign each output pin to a physical device pin using \textbf{Assignment} $\to$ \textbf{Pin Planner}, per Table~\ref{tab:exp3-pins} (refer to the DE2 user manual for pin locations).
  \item Save, then recompile the project.
  \item Download the configuration to the DE2 board via \textbf{Tools} $\to$ \textbf{Programmer} $\to$ \textbf{Hardware Setup} $\to$ USB-Blaster $\to$ \textbf{Start}.
  \item Record the pattern displayed on the seven-segment display.
  \item Change the assignment to \verb|assign q[6:0] = 7'b1111000;|, save, recompile, and download again. Record the new displayed pattern.
\end{enumerate}

\begin{codepanel}{counta.v}{verilogstyle}
\begin{lstlisting}
module counta(q);
  output [6:0] q;
  assign q[6:0] = 7'b0010000;
endmodule
\end{lstlisting}
\end{codepanel}
\vspace{-4pt}
\begin{center}\small\color{slategray}Listing \refstepcounter{equation}\theequation\label{lst:exp3-counta}: \texttt{counta.v} --- static seven-segment pattern.\end{center}

\begin{boxedtable}
\centering
\begin{tabular}{clcl}
\toprule
\# & Node name & Direction & Pin \\
\midrule
1 & Q[6] & Output & PIN\_V13 \\
2 & Q[5] & Output & PIN\_V14 \\
3 & Q[4] & Output & PIN\_AE11 \\
4 & Q[3] & Output & PIN\_AD11 \\
5 & Q[2] & Output & PIN\_AC12 \\
6 & Q[1] & Output & PIN\_AB12 \\
7 & Q[0] & Output & PIN\_AF10 \\
\bottomrule
\end{tabular}
\captionof{table}{Seven-segment output pin assignment (DE2 board).}
\label{tab:exp3-pins}
\end{boxedtable}

\end{procedure}
